% Part:sets-functions-relations
% Chapter: size-of-sets
% Section: enumerations-alt

\documentclass[../../../include/open-logic-section]{subfiles}

\begin{document}

\olfileid{sfr}{siz}{enm-alt}

\olsection{లెక్కింపులు, \usetoken{S}{enumerable} సమితులు}

\begin{editorial}
  సమితి సిద్ధాంతంలో చేసినట్లుగా, ఈ విభాగం లెక్కింపులను $\Nat$తో
  (లేదా దాని ఆరంభ ఖండాలతో) ద్విగుణ ప్రమేయాలుగా నిర్వచిస్తుంది.
  అందువల్ల ఇది \olref[enm]{sec}లోని నిర్వచనాలతో కొద్దిగా విభేదించి,
  అక్కడి ఉదాహరణలన్నిటినీ పునరావృతం చేస్తుంది. ఆ విభాగం కంటే ఇది
  కొంచెం సంక్షిప్తం కూడా.
\end{editorial}

ఒక పరిమిత సమితిని దాని !!{element}sను వరుసగా పేర్కొనడం ద్వారానే
తెలియజేయవచ్చు. సమితిని ఇలా నిర్వచించినప్పుడు మనం ఇదే చేస్తాం:
\[
  A = \{a_1, a_2, \ldots, a_n\}.
\]
!!{element}s $a_1$, \dots, $a_n$ అన్నీ వేర్వేరని అనుకుంటే, ఇది $A$కు,
మొదటి $n$ సహజ సంఖ్యలు $0$, \dots, $n-1$కు మధ్య !!a{bijection}ను
ఇస్తుంది. తిరిగి, ప్రతి పరిమిత సమితిలో పరిమిత సంఖ్యలోనే !!{element}s
ఉంటాయి కనుక, ప్రతి పరిమిత సమితినీ అటువంటి సహసంబంధంలో పెట్టవచ్చు.
మరోలా చెప్పాలంటే, $A$ పరిమితమైతే, $A$కు, $\{0, \dots, n-1\}$కు
మధ్య !!a{bijection} ఉంటుంది; ఇక్కడ $n$ అనేది $A$లోని !!{element}s
సంఖ్య.

కొన్ని రకాల అనంత సమితులను అనుమతిస్తే, కొన్ని అనంత సమితులను కూడా
లెక్కించడానికి అనుమతిస్తాం. ఒక అనంత సమితికి, $\Nat$ మొత్తానికి మధ్య
!!a{bijection} ఉంటే ఆ సమితి లెక్కించబడుతుందని చెప్పి దీన్ని కచ్చితంగా
వ్యక్తం చేయవచ్చు.

\begin{defn}[లెక్కింపు, సమితి-సిద్ధాంత రూపం]
$A$ సమితి యొక్క \emph{లెక్కింపు} అంటే, దాని పరిధి $A$గానూ, ప్రవేశం
సహజ సంఖ్యల ఆరంభ సమితి $\{0, 1, \ldots, n\}$ లేదా సహజ సంఖ్యల
మొత్తం సమితి $\Nat$గానూ ఉన్న !!a{bijection}.
\end{defn}

\begin{explain}
\emph{లెక్కింపు} అనే పదాన్ని ఇలా వాడటానికి సహజమైన ఆధారం ఉంది. $A$
సమితిని లెక్కించామని చెప్పడం అంటే, $A$లోని మూలకాలను వరుసగా లెక్కించడానికి
ఉపయోగపడే !!a{bijection} $f$ ఉందని చెప్పడమే. $0$వ మూలకం $f(0)$,
మొదటిది $f(1)$, \ldots, $n$వది $f(n)$, \ldots.
\footnote{అవును, మనం $0$ నుంచి లెక్కిస్తాం. కావాలంటే $1$తో కూడా
మొదలుపెట్టవచ్చు. దానివల్ల పెద్ద తేడా ఏమీ ఉండదు; $\Nat$కు బదులు
$\PosInt$ను పెట్టాల్సి వస్తుందంతే.} కింది నిర్వచనాన్ని చేరిస్తే
దీనికి కారణం ఇంకా స్పష్టమవుతుంది:
\end{explain}

\begin{defn}
  \ollabel{defn:enumerable}
  $A = \emptyset$ లేదా $A$కు ఒక లెక్కింపు ఉంటేనే $A$ సమితి
  !!{enumerable}. $A$ !!{enumerable} కాకపోతేనే $A$
  !!{nonenumerable} అంటాం.
\end{defn}

\begin{explain}
అందువల్ల ఒక సమితి !!{enumerable} కావడానికి అవసరమైన మరియు సరిపడే
షరతు అది ఖాళీ సమితి కావడం లేదా దాని !!{element}sను ఒక లెక్కింపుతో
వరుసగా లెక్కించగలగడం.
\end{explain}

\begin{ex}
$\Id{\Nat}(n) = n$గా ఇచ్చే తాదాత్మ్య ప్రమేయం
$\Id{\Nat} \colon \Nat \to \Nat$, సహజ సంఖ్యలను లెక్కించే ప్రమేయం.
\emph{ధన} సహజ సంఖ్యలు $\Nat^+ = \Nat \setminus \{0\}$ను లెక్కించే
ప్రమేయం $g(n) = n + 1$; అంటే, ఉత్తరవర్తి ప్రమేయం.
\end{ex}

\begin{prob}
$A$ సమితి !!{enumerable} కావడానికి అవసరమైన మరియు సరిపడే షరతు
$A = \emptyset$ లేదా !!a{surjection} $f\colon \Nat \to A$ ఉండటం అని
చూపండి. $A$ !!{enumerable} కావడానికి అవసరమైన మరియు సరిపడే షరతు
!!a{injection} $g\colon A \to \Nat$ ఉండటం అని చూపండి.
\end{prob}

\begin{ex}
ఈ విధంగా ఇచ్చిన $f\colon \Nat \to \Nat$, $g \colon \Nat \to \Nat$
ప్రమేయాలు
\begin{align*}
  f(n) & = 2n \text{ మరియు}\\
  g(n) & = 2n+1
\end{align*}
వరుసగా సరి సహజ సంఖ్యలను, బేసి సహజ సంఖ్యలను లెక్కిస్తాయి. కానీ వాటిలో
ఏదీ !!{surjective} కాదు; అందువల్ల వాటిలో ఏదీ $\Nat$ యొక్క లెక్కింపు
కాదు.
\end{ex}

\begin{prob}
వర్గ సంఖ్యలు $1$, $4$, $9$, $16$, \dots{}కు ఒక లెక్కింపును నిర్వచించండి.
\end{prob}

\begin{ex}
$x$ను దాని కంటే తక్కువ కాని సమీప పూర్ణాంకానికి పెంచి సవరించే
\emph{సీలింగ్} ప్రమేయం $\lceil x \rceil$ అనుకుందాం. అప్పుడు ఇలా
ఇచ్చిన $f \colon \Nat \to \Int$:
\[
  f(n) = (-1)^{n} \left\lceil\tfrac{n}{2}\right\rceil
\]
పూర్ణాంకాల సమితి $\Int$ను ఈ విధంగా లెక్కిస్తుంది:
\[
\begin{array}{c c c c c c c c}
f(0) & f(1) & f(2) & f(3) & f(4) & f(5) & f(6) & \dots \\ \\
\big\lceil \tfrac{0}{2} \big\rceil & -\big\lceil \tfrac{1}{2}\big\rceil &  \big\lceil \tfrac{2}{2} \big\rceil & -\big\lceil \tfrac{3}{2} \big\rceil & \big\lceil \tfrac{4}{2} \big\rceil  & -\big\lceil \tfrac{5}{2}\big\rceil & \big\lceil \tfrac{6}{2} \big\rceil & \dots \\ \\
0 & -1 & 1 & -2 & 2 & -3 & 3& \dots
\end{array}
\]
ధన, ఋణ పూర్ణాంకాల మధ్య అటూ ఇటూ ``గెంతుతూ'' $f$, $\Int$ విలువలను
ఎలా ఉత్పత్తి చేస్తుందో గమనించండి. $f$ను సందర్భాలవారీగా ఇలా
నిర్వచించిందని కూడా భావించవచ్చు:
\[
f(n) = \begin{cases}
  \frac{n}{2} & \text{$n$ సరి అయితే}\\
  -\frac{n+1}{2} & \text{$n$ బేసి అయితే}
  \end{cases}
\]
\end{ex}

\begin{prob}
$A$, $B$ !!{enumerable} అయితే $A \cup B$ కూడా అలాంటిదే అని చూపండి.
\end{prob}

\begin{prob}
$A_1$, $A_2$, \dots, $A_n$ అన్నీ !!{enumerable} అయితే
$A_1 \cup \dots \cup A_n$ కూడా అలాంటిదే అని $n$పై గణిత ఆగమనం
ద్వారా చూపండి.
\end{prob}

\end{document}
